\chapter{Aboriginal and Torres Strait Islander Health}

\section*{Definition and Overview}
For Aboriginal and Torres Strait Islander peoples, health is understood far more broadly than the absence of disease. The holistic concept of health embraces physical, social, emotional, and cultural wellbeing, together with connection to Country, community, culture, family, and ancestors. It encompasses spiritual wellbeing and the collective wellbeing of the community, not just the individual. Effective health care therefore requires cultural safety, genuine respect for identity and community, and the delivery of services in genuine partnership with Aboriginal and Torres Strait Islander people. This chapter outlines the main health concerns affecting these communities, the social and historical factors that shape them, and the approaches that improve outcomes.

\section*{Epidemiology}
Aboriginal and Torres Strait Islander peoples experience a substantially higher burden of illness than the wider Australian population. Life expectancy is approximately 8--10 years shorter than for non-Indigenous Australians. Rates are elevated for cardiovascular disease, type 2 diabetes, chronic kidney disease, chronic lung disease, some cancers, and ear disease, and psychological distress, mental ill-health, and suicide are also more common. Several serious infections that are now rare in the wider population, including rheumatic heart disease and tuberculosis, continue to affect these communities. There have been meaningful improvements in recent decades, including falling infant mortality and declining smoking rates, yet significant gaps persist.

\section*{Aetiology and Pathophysiology}
The causes of the health gap are social, cultural, and historical rather than biological.

\begin{itemize}
    \item \textbf{Social determinants:} lower average income, educational attainment, employment, and housing standards than the wider population
    \item \textbf{The legacy of colonisation:} dispossession of land, forced removal of children, and the intergenerational trauma that follows
    \item \textbf{Racism and discrimination:} experienced in daily life and within health services themselves, acting as a barrier to care
    \item \textbf{Barriers to access:} distance, cost, and a shortage of culturally safe services, particularly in rural and remote areas
    \item \textbf{Higher rates of smoking} and some dietary, alcohol, and activity patterns
    \item \textbf{Environmental factors:} overcrowded and poor-quality housing, limited water security, and restricted access to affordable, healthy food
    \item \textbf{Protective strengths:} strong families, communities, cultural identity, and resilience, which are powerful assets for health
\end{itemize}

\section*{Clinical Features}
People may present in a range of ways that reflect the burden of chronic and infectious disease in the community.

\begin{itemize}
    \item Symptoms of chronic conditions such as diabetes, heart disease, or kidney disease, which may already be advanced when first detected
    \item Repeated infections, including ear, chest, and skin infections, especially in children
    \item Mental distress that may present as physical symptoms, irritability, or social withdrawal
    \item Complications such as kidney failure, blindness, or lower-limb amputations when disease has been detected late
    \item Presentations in which physical, emotional, spiritual, family, and community concerns are intertwined
\end{itemize}

\section*{Differential Diagnosis}
In any presentation, clinicians should consider the full range of physical, mental, and social causes.

\begin{itemize}
    \item Physical illness versus emotional, social, or cultural distress
    \item A common chronic disease versus other causes of the same symptoms
    \item A single condition versus several coexisting conditions
    \item The effects of current medicines or previous treatments
    \item Family and community factors that may be causing, contributing to, or worsening symptoms
\end{itemize}

\section*{When to Seek Medical Care}
Aboriginal and Torres Strait Islander adults are encouraged to have a regular annual health check, funded through Medicare, even when feeling well, so that risk factors and early disease can be identified. Care should be sought sooner for any new, persistent, or worrying symptom such as chest pain, breathlessness, fatigue, unexplained weight change, or mental distress.

\textbf{Call 000 (or local emergency services) for:}
\begin{itemize}
    \item Chest pain, difficulty breathing, or collapse
    \item Sudden weakness, facial drooping, or difficulty speaking, which may indicate a stroke
    \item Severe abdominal pain or heavy bleeding
    \item Confusion, severe agitation, or a mental health crisis with a risk of harm
    \item Any sudden, severe pain or other emergency you are worried about
\end{itemize}

\section*{Diagnosis and Investigations}
Care is organised around regular preventive assessment and the early detection of common chronic conditions.

\begin{itemize}
    \item \textbf{Annual health check:} the Medicare-funded health assessment (item 715), which reviews risk factors, symptoms, and social and emotional wellbeing
    \item \textbf{Blood pressure, weight, and body mass index} measured at each visit
    \item \textbf{Blood tests:} blood glucose or HbA1c for diabetes, lipid profile, kidney function, and liver function
    \item \textbf{Cardiovascular risk assessment} and screening for rheumatic heart disease and other important infections
    \item \textbf{Ear and hearing checks} in children, and regular eye checks for diabetic retinopathy
    \item \textbf{Screening for mental health and social and emotional wellbeing}, conducted in a culturally safe manner
    \item \textbf{Involvement of family and community} where the person wishes, since decisions are often made collectively
\end{itemize}

\section*{Management}
Effective management is prevention-focused, team-based, and delivered in culturally safe ways.

\begin{itemize}
    \item \textbf{Culturally safe primary care,} often through Aboriginal Community Controlled Health Organisations (ACCHOs) staffed by Aboriginal and Torres Strait Islander health workers
    \item \textbf{Chronic disease management} with a multidisciplinary team including doctors, nurses, allied health professionals, and community health workers
    \item \textbf{Prevention-focused care} that treats infections such as strep throat and skin sores early to prevent complications including rheumatic heart disease and kidney disease
    \item \textbf{Coordinated care} across hospital, community, and outreach services, particularly for rural and remote communities
    \item \textbf{Care that involves family and community} and respects cultural obligations and priorities
    \item \textbf{Culturally adapted health education} around smoking, nutrition, physical activity, and diabetes
    \item \textbf{Regular follow-up} to maintain control of chronic conditions and detect complications early
\end{itemize}

\section*{Complications}
\begin{itemize}
    \item Late diagnosis of cancer, heart disease, and other conditions
    \item Kidney failure requiring long-term dialysis or transplantation
    \item Blindness and foot problems, including amputations, from poorly controlled diabetes
    \item Rheumatic heart disease following untreated streptococcal throat infections
    \item Permanent hearing loss from recurrent ear infections in childhood
    \item Mental health crises, self-harm, and suicide
\end{itemize}

\section*{Prognosis and Outcomes}
Outcomes have improved in several areas over recent decades, including infant mortality, smoking rates, and the earlier detection of some cancers. With early detection and culturally safe, well-coordinated care, chronic conditions can be managed so that people live long, full, and healthy lives. Sustained progress also depends on action on the social determinants of health, including housing, education, employment, and income. Evidence increasingly shows that health outcomes are best when services are led by, and delivered in partnership with, Aboriginal and Torres Strait Islander communities.

\section*{Prevention and Self-Care}
\begin{itemize}
    \item Attend an annual health check, even when feeling well
    \item Quit smoking, with support from a GP or Aboriginal health service
    \item Eat well, stay active, and keep alcohol within recommended limits
    \item Take prescribed medicines for conditions such as diabetes and high blood pressure
    \item Keep children's vaccinations and ear, hearing, and general health checks up to date
    \item Treat sore throats and skin sores promptly to prevent serious complications
    \item Stay connected with family, community, and culture, which strongly support wellbeing
    \item Seek help early for mental health concerns, for yourself or someone close to you
\end{itemize}

\section*{Sources and Further Reading}
The following organisations provide further information on Aboriginal and Torres Strait Islander health for students, clinicians, and community members.

\begin{itemize}
    \item Australian Government Department of Health. \textit{Aboriginal and Torres Strait Islander health programs and policy.} https://www.health.gov.au/
    \item Australian Institute of Health and Welfare. \textit{Aboriginal and Torres Strait Islander health data and reports.} https://www.aihw.gov.au/
    \item NACCHO. \textit{National Aboriginal Community Controlled Health Organisation.} https://www.naccho.org.au/
    \item World Health Organization. \textit{Indigenous health and social determinants.} https://www.who.int/
    \item Centers for Disease Control and Prevention. \textit{Social determinants of health.} https://www.cdc.gov/
    \item MedlinePlus. \textit{Social determinants of health.} https://medlineplus.gov/healthdisparities.html
\end{itemize}
